\section{Experiments}

For fair and reproducible comparison with other models, we used the LJSpeech dataset~\citep{ljspeech17} in which many speech synthesis models are trained. LJSpeech consists of 13,100 short audio clips of a single speaker with a total length of approximately 24 hours. The audio format is 16-bit PCM with a sample rate of 22 kHz; it was used without any manipulation. HiFi-GAN was compared against the best publicly available models: the popular mixture of logistics (MoL) WaveNet~\citep{oord2018parallel} implementation~\citep{Yamamoto18wavenet}
\footnote{In the original paper~\citep{oord2018parallel}, an input condition called linguistic features was used, but this implementation uses mel-spectrogram as the input condition as in the later paper~\citep{shen2018natural}.} and the official implementation of WaveGlow~\citep{WaveGlowRepo} and MelGAN~\citep{MelGANRepo}. We used the provided pretrained weights for all the models.

To evaluate the generality of HiFi-GAN to the mel-spectrogram inversion of unseen speakers, we used the VCTK multi-speaker dataset~\citep{veaux2017cstr}, which consists of approximately 44,200 short audio clips uttered by 109 native English speakers with various accents. The total length of the audio clips is approximately 44 hours. The audio format is 16-bit PCM with a sample rate of 44 kHz. We reduced the sample rate to 22 kHz. We randomly selected nine speakers and excluded all their audio clips from the training set. We then trained MoL WaveNet, WaveGlow, and MelGAN with the same data settings; all the models were trained until 2.5M steps.

To evaluate the audio quality, we crowd-sourced 5-scale MOS tests via Amazon Mechanical Turk. The MOS scores were recorded with 95\% confidence intervals (CI). Raters listened to the test samples randomly, where they were allowed to evaluate each audio sample once. All audio clips were normalized to prevent the influence of audio volume differences on the raters. 
All quality assessments in Section 4 were conducted in this manner, and were not sourced from other papers.

The synthesis speed was measured on GPU and CPU environments according to the recent research trends regarding the efficiency of neural networks~\citep{Kumar2018melgan, zhai2020squeezewave, tan2020efficientdet}. The devices used are a single NVIDIA V100 GPU and a MacBook Pro laptop (Intel i7 CPU 2.6GHz). Additionally, we used 32-bit floating point operations for all the models without any optimization methods.

To confirm the trade-off between synthesis efficiency and sample quality, we conducted experiments based on the three variations of the generator, $V1$, $V2$, and $V3$ while maintaining the same discriminator configuration. For $V1$, we set $h_u=512$, $k_u=[16, 16, 4, 4]$, $k_r=[3,7,11]$, and $D_r=[[1,1],[3,1],[5,1]] \times 3]$. $V2$ is simply a smaller version of $V1$, which has a smaller hidden dimension $h_u=128$ but with exactly the same receptive fields. To further reduce the number of layers while maintaining receptive fields wide, the kernel sizes and dilation rates of $V3$ were selected carefully. The detailed configurations of the models are listed in Appendix \hyperref[app-a1]{A.1}. We used 80 bands mel-spectrograms as input conditions. The FFT, window, and hop size were set to 1024, 1024, and 256, respectively. The networks were trained using the AdamW optimizer~\citep{loshchilov2017decoupled} with $\beta_1=0.8$, $\beta_2=0.99$, and weight decay $\lambda=0.01$. The learning rate decay was scheduled by a $0.999$ factor in every epoch with an initial learning rate of $2\times 10^{-4}$.
